% RESEARCH_DRAFT_SHA256: f68525b7e177116766fbc2ab32842511e75f9c66dcbcab9835b48bdf1ad4adee
% Layout and source-reference wrappers expanded for the web renderer.
% FULL_SOLUTION_SCOPE: {"target": "Question47: do the interpolated generators always have zero set equal to the finite rational design under every stated hypothesis?", "result": "No. The displayed four-point rational design satisfies every hypothesis, but both interpolated minimal-generator equations vanish on the entire diagonal.", "coverage": "A universal affirmative assertion is disproved by one fully verified admissible rational design; the extra point (3,3) is rational and also a complex zero.", "dependencies": "All independence, interpolation and zero-set calculations are proved in the manuscript; no external theorem is a proof dependency.", "human_verification": "None; model-assisted self-audit and Codex GPT-6 mathematical review only.", "unverified": "No independent human or formal proof verification.", "novelty": "Negative answer already recorded by AIM and attributed to Bernd Sturmfels; the witness is rederived here without novelty or historical-coordinate attribution."}
\documentclass[11pt]{article}
\usepackage[margin=25mm]{geometry}
\usepackage{fontspec,unicode-math}
\setmainfont{Latin Modern Roman}
\setmathfont{Latin Modern Math}
\usepackage{amsmath,amssymb,mathtools,xurl,hyperref}
\hypersetup{colorlinks=true,urlcolor=blue,linkcolor=blue}
\setcounter{secnumdepth}{0}
\setlength{\parindent}{0pt}
\setlength{\parskip}{5pt}
\setlength{\emergencystretch}{3em}
\newcommand{\sref}[2]{\hyperlink{source:#1:#2}{[S#2]}}
\newcommand{\cataloguescope}[1]{\paragraph{Source record.}#1}
\newcommand{\floor}[1]{\left\lfloor#1\right\rfloor}
\begin{document}
\section*{Research-attempt provenance}
This standalone research attempt retains the complete original definition body below,
including its historical source and preparation notes. Those notes describe the
earlier catalogue preparation; the new attempt and its verification are identified
separately after that body. No original ID, catalogue status, or source statement is changed.

\section{A four-point counterexample and the border-basis repair}\label{statement}

\subsection{Definitions and mathematical statement}

\paragraph{Shared notation.}Standard number systems are $\mathbb N,\mathbb Z,\mathbb Q,\mathbb R,\mathbb C$. Specialized notation and index ranges follow the cited source and the definitions below.


For a finite rational design $D$, its design ideal is $I(D)=\{f:f(v)=0\text{ for all }v\in D\}$. Monomials are independent modulo $I(D)$ when their residue classes are linearly independent in the quotient ring. Standard monomials of a monomial ideal are those not in that ideal. Interpolating a generator $m_i$ in a quotient basis yields $f_i=m_i-h_i\in I(D)$; the question asks whether their common zero set introduces extra points.
\paragraph{Record scope and qualifications.}
The catalogue title is retained as metadata. Named certificates, contracts, bounds or frameworks in that title are not asserted as results of the original source extract.
The source uses inconsistent dimension letters, which are retained literally. The next canonical record reports a negative workshop answer; no affirmative zero-set theorem is claimed. The supplied catalogue marks this record solved; its present research status was not independently checked.
The source text below is rendered literally where necessary; damaged formulas and composite extracts are not silently repaired into a new mathematical claim.
\paragraph{Statement recorded in the supplied source.}
Question 47. [Pistone] Take a design of N distinct rational vectors D = \{v1,..., v N \} ⊆ Qd

and xα1,..., x αN are monomials in Q[x1,..., x n] which are linearly independent modulo the design ideal I(D). Further suppose that these monomials are equal to the set of standard monomials of some monomial ideal M.Suppose that M = 〈m1,..., m r〉. Since the monomials xα1,..., x αN are a basis for a quotient space, there is a unique representation of the mi as linear combinations of the xαj.This produces r polynomials f1,..., f r which belong to I(D). Question: are the only points which vanish on these polynomials the design points; that is, does V (f1,..., f r) = D? \href{https://aimath.org/WWN/compalgstat/compalgstat.pdf}{[S2]}
\paragraph{Source status.}
The supplied catalogue labels this source record solved; that metadata was checked against the supplied snapshot on 2026-10-11. The label is preserved for provenance and is not an independent verification of present open status or mathematical correctness. This is documentation of a source-labelled solved record, not a newly proposed unsolved problem.
\subsection{Short English statement}
The historical question asks whether the constructed interpolation equations cut out exactly the design points.
\subsection{Sources}
\begin{itemize}
\item[S1] Ulam AI and the UnsolvedMath Contributors, UnsolvedMath: A Curated Collection of Open Mathematics Problems, catalogue entry. Catalogue curation and metadata: CC BY 4.0. The source extract is retained; definitions, English summaries and qualifications are editorial additions.
\url{https://huggingface.co/datasets/ulamai/UnsolvedMath}.
\item[S2] AIM Workshop Problem Lists, source linked by the catalogue.
\url{https://aimath.org/WWN/compalgstat/compalgstat.pdf}.
\end{itemize}
\paragraph{Preparation and review.}
The original source and UnsolvedMath contributors retain attribution for the source extract. Codex (GPT-6; exact serving revision not exposed) assisted with definitions, English summaries, scope notes and formatting. No new human mathematical authorship or independently verified solution is asserted. Catalogue identity checked against the complete supplied registry and current repository. Distinct upstream records, including repeated questions and broad or damaged extracts, are retained for complete coverage of the source CS category. A repeated record is not a new independent conjecture. Complete mathematical-equivalence and current-literature review has not been performed. Related canonical identities: 20001224. Preparation and final review follow contribution prompts C and D, with source gaps disclosed. Structural and rendering checks do not establish a new theorem or open status.
\label{end:statement}

\section{Literature and context check}
The primary AIM document, Question 47 and its following sentence, explicitly records a negative workshop answer attributed to Bernd Sturmfels. The source does not supply coordinates. The calculation below reconstructs a witness; neither its novelty nor its identity with Sturmfels' historical example is asserted. The catalogue's solved label is preserved.
\section{Research attempt: A complete witness for the already recorded negative answer}
Work in $\mathbb Q[x,y]$, and take
\[
D=\{(0,0),(1,1),(2,2),(0,1)\},\qquad
\mathcal O=\{1,x,y,xy\},\qquad M=(x^2,y^2).
\]
These are four distinct rational points, and $\mathcal O$ is exactly the set of standard monomials of $M$: a monomial avoids $M$ precisely when both exponents are at most one. To prove independence modulo $I(D)$, suppose $a+bx+cy+dxy$ vanishes on $D$. On the first three points the polynomial $a+(b+c)t+dt^2$ vanishes at $0,1,2$. Substitution gives $a=0$, $b+c+d=0$, and $2(b+c)+4d=0$, whence $d=0$ and $b+c=0$. Substitution at $(0,1)$ gives $c=0$, hence $b=0$. Thus the evaluation map from the four-dimensional span of $\mathcal O$ to $\mathbb Q^D$ is bijective. Every polynomial has a unique interpolating representative in that span, because its four values do.

The two minimal generators interpolate as
\[
x^2\equiv xy,\qquad y^2\equiv xy-x+y\pmod{I(D)}.
\]
Indeed these identities hold on the diagonal and at $(0,1)$. The source construction therefore gives
\[
f_1=x^2-xy=x(x-y),\qquad
f_2=y^2-xy+x-y=(y-x)(y-1).
\]
Their zero set over any characteristic-zero field is the entire diagonal together with $(0,1)$. To prove equality, points with $x=y$ satisfy both equations. At a common zero with $x\ne y$, the first equation forces $x=0$ and the second forces $y=1$. In particular $(3,3)$ is an extra rational point. Every hypothesis of the precise historical question holds, so its proposed affirmative statement fails.

The attempted inference that independence of the proposed standard monomials implies that the interpolated corner relations generate $I(D)$ fails here. All corner equations share $x-y$ as a factor, even though their evaluations impose four independent interpolation conditions. An explicit additional relation removes the excess component:
\[
h=xy^2-3xy+2x=x(y-1)(y-2).
\]
It vanishes on $D$. On the excess diagonal it restricts to $x(x-1)(x-2)$, and it also vanishes at $(0,1)$. Consequently $V(f_1,f_2,h)=D$. This is an instance-specific repair, not a proof that arbitrary corner relations suffice.
\section{Verification and adversarial audit}
The four interpolation conditions, the minimal-generator requirement, and the extra rational zero were checked directly in the displayed equations. The proof classifies every common zero, rather than relying on plotting or numerical solving. The added relation proves a zero-set equality only; no unproved equality of ideals or multiplicities is used.
\paragraph{Audit corrections.} A negative answer to the old question is explicitly distinguished from a newly solved open problem. The source's attribution is retained, but the coordinates are not attributed to Sturmfels without evidence.
\paragraph{External computation.} None; symbolic arguments and boundary checks only.
\section{Scope and remaining gaps}
The original universal yes/no question is completely answered in the negative by
the displayed admissible design and extra rational zero. The negative answer was
already recorded in the primary source; no new resolution is claimed. The added
relation is an auxiliary repair of this example, not a general border-basis theorem.
\section{Full-solution scope declaration}
The full original target asks whether the specified construction always cuts out
exactly the finite rational design. The four-point example disproves that assertion
with all its hypotheses intact. The proof explicitly checks distinct rational
points, the standard-monomial condition, quotient independence, both unique
interpolants, and an extra rational common zero. All mathematical dependencies
are proved above. Codex GPT-6 supplied model-assisted mathematical review; no
independent human or formal proof verification occurred. The negative answer is
known from the cited source, and no novelty or identification of these coordinates
with the historical example is asserted. General repair theorems and efficiency
questions are not part of the original yes/no target and are not claimed solved.
\paragraph{Integration review.} Codex (GPT-6; exact serving revision not exposed) reviewed the written arguments and expanded layout and source-reference wrappers. Repository metadata and the production web renderer were checked using Python 3.13 and JavaScript with MathJax 3.2.2. This model-assisted review is not independent human proof verification; no PDF compilation or formal proof checking is asserted.

\section{Completion Estimate}
\noindent\textbf{COMPLETION: 100\%}
Every hypothesis of the original universal assertion is met by the counterexample,
and its conclusion fails at the explicit rational point $(3,3)$. This measures
complete original-target coverage, not novelty.
\section{References}
The original source references and catalogue attribution are preserved in the
Sources section above. The following URLs were actually consulted for this attempt;
retention of an additional original reference does not claim fresh consultation of it.
\begin{itemize}
\item \url{https://aimath.org/WWN/compalgstat/compalgstat.pdf} (consulted 2026-10-11). 
\end{itemize}
\section{Model and contributor disclosure}
OpenAI gpt-6-astra, requested reasoning configuration ultra, generated this research
attempt on 2026-10-11. The exact serving revision is not exposed. Codex GPT-6
performed orchestration, repository-format review, and a model-assisted mathematical
audit of the written arguments; its exact serving revision
is also not exposed. The mathematical route, boundary checks and final adversarial
self-audit followed contribution prompts A then B in
\texttt{docs/contribution\_prompts.md}, after reading \texttt{CONTRIBUTING.md}.
Python 3.14.2 assembled the standalone documents; any mathematical computation is
specified in Verification. No theorem prover was used. No human mathematical
authorship, substantial human proof correction, or independent human proof checking
is asserted. Model-assisted review, source retrieval, and structural checks do not
establish correctness or novelty. Full private reasoning transcripts are unavailable
as reproducibility artifacts; the complete arguments, inputs and outcomes actually
used are stated here or in the companion batch audit.
\section{Final status}
\textbf{Solved as a rederivation of a known negative answer.} The counterexample
completely resolves the original universal zero-set question. No new theorem or
historical attribution of the chosen coordinates is claimed.
\end{document}
